\chapter{Literature Review}
\label{ch:lit}

This project sits where several research areas meet: classical timing and cache attacks on AES, statistical leakage assessment, constant-time implementation techniques, microarchitectural attacks on Apple processors, and the recent work on how compilers break constant-time code. The work is grouped by theme. For each theme, the review notes what this project borrows and what remains open.

\section{Timing and cache attacks on AES}
\label{sec:lit_attacks}

Kocher~\cite{kocher96} introduced timing attacks in 1996, recovering private keys from Diffie-Hellman, RSA and DSS implementations whose running time depended on secret exponent bits. The paper established that a mathematically secure algorithm can be broken through a physical property of its implementation.

Bernstein~\cite{bernstein05} applied the idea to AES in 2005. He showed that the time taken by OpenSSL's table-based AES depended on which S-box table entries were already in the cache, and he recovered a key from a server that reported its own timing. His report also listed obstacles to writing constant-time table lookups on real processors. Osvik, Shamir and Tromer~\cite{osvik06} systematised cache attacks on AES with two techniques that are still the basis of the field: \emph{Evict+Time}, where the attacker evicts cache sets and measures the victim's total encryption time, and \emph{Prime+Probe}, where the attacker fills cache sets and later measures which of its own lines were evicted.

\textbf{Borrowed for this project:} the cache-eviction measurement condition of Part I is a simplified Evict+Time setting, and the T-table implementation reproduces the leaky design these papers attacked. \textbf{Open point:} these attacks were developed on x86 processors with 64-byte cache lines and small L1 caches. How they appear on a processor with 128-byte lines and a 128~KB L1 data cache (Apple M4) is an empirical question that Part I measures.

\section{Statistical leakage assessment}
\label{sec:lit_tvla}

Goodwill, Jun, Jaffe and Rohatgi~\cite{tvla} proposed Test Vector Leakage Assessment (TVLA) at the NIST Non-Invasive Attack Testing workshop in 2011. Instead of mounting a specific attack, TVLA compares measurements taken with a fixed input against measurements with random inputs, using Welch's $t$-test~\cite{welch47} with a threshold of $\lvert t\rvert = 4.5$. The method was designed for power and electromagnetic traces but transfers directly to timing. Reparaz, Balasch and Verbauwhede~\cite{dudect} did exactly that with \emph{dudect} (``Dude, is my code constant time?'', DATE 2017): many timing measurements of a function under fixed and random input classes, randomised interleaving of classes, cropping of outliers at several percentiles, and a running Welch's $t$. The \code{dudect-bencher} crate~\cite{dudectbencher} implements the same method for Rust.

\textbf{Borrowed:} the fixed-versus-random protocol, the 4.5 threshold, randomised interleaving and percentile cropping. \textbf{Added in Part I:} explicit positive and negative controls, robust rank-based statistics, Holm correction across all tests, and a minimum detectable difference that bounds every non-detection.

\section{Constant-time AES implementations}
\label{sec:lit_ctaes}

Käsper and Schwabe~\cite{kasper09} showed that AES can be implemented without any secret-indexed memory access through \emph{bit-slicing}, processing eight blocks in parallel with bitwise operations, while remaining fast on processors with wide registers. Hamburg~\cite{hamburg09} reached the same goal with vector permute instructions that perform a 16-entry table lookup inside a register, so no memory address depends on the secret. Both approaches are used in OpenSSL (\code{bsaes} and \code{vpaes}). Hardware AES instructions remove the problem altogether, because the S-box is implemented inside the execution unit.

\textbf{Borrowed:} the constant-time implementation of Part I follows the same principle by computing the S-box as $A(a^{254})$ in a SWAR fashion, without tables or branches. \textbf{Open point:} these designs argue constant-time behaviour at the source and instruction level; on a specific modern processor, the claim should also be tested empirically.

\section{Microarchitectural attacks on Apple Silicon}
\label{sec:lit_apple}

Three recent attacks show that Apple's processors contain optimisation hardware that can undermine source-level constant-time reasoning. GoFetch~\cite{gofetch} (USENIX Security 2024) exploited the \emph{data memory-dependent prefetcher} (DMP) of Apple M-series chips: if a value loaded from memory looks like a pointer, the DMP may prefetch it, so data values, not just addresses, influence the cache. The authors extracted keys from constant-time implementations of classical and post-quantum cryptography. SLAP~\cite{slap} (IEEE S\&P 2025) and FLOP~\cite{flop} (USENIX Security 2025) showed that Apple processors include a \emph{load address predictor} and a \emph{load value predictor}, and that their mispredictions can be used to read memory speculatively. FLOP reports that the load value predictor activates for loads that repeatedly return the same value.

\textbf{Relevance:} Part I runs on an Apple M4. One of its results, branch-free code flagged only when the same input is repeated in long batches, is consistent in spirit with predictor state that persists across executions, although Part I does not confirm a mechanism.

\section{Compiler-induced constant-time violations}
\label{sec:lit_compilers}

The most directly relevant work for Part II asks whether compilers break code that was written to be constant-time.

\textbf{Breaking Bad} (Schneider, Lain, Puddu, Dutly and Čapkun, ASIA CCS 2025)~\cite{breakingbad} is a large-scale study of compiler-induced secret-dependent operations. It analyses eight cryptographic libraries across six processor architectures with many compiler versions and optimisation settings, and finds compiler-induced secret-dependent operations in some of the most highly regarded hardened libraries. The authors conclude that defensive source-level programming is not sufficient once compiler optimisations are applied. The evaluated libraries are C libraries; the paper does not identify Rust crates among its targets.

\textbf{Fun with Flags} (Geimer and Maurice, PriSC 2026)~\cite{funwithflags} builds a dataset of compiler-introduced constant-time violations in GCC and LLVM and traces most of them to a small number of optimisation passes. It proposes disabling those passes through compiler flags, which removes the compiler-introduced violations in its dataset at low performance cost.

\textbf{DOCC} (Gerlach, Pietsch and Schwarz, FC 2025)~\cite{docc} is an automated binary testing pipeline that compiles cryptographic code under many compiler configurations and tests the binaries for data-obliviousness violations. It reports violations in 93 of 127 tested algorithms.

\textbf{RoboCop} (Kolosick et al., PLDI 2025)~\cite{robocop} approaches the problem from the defensive side: a compiler-based toolchain that generates protected cryptographic library code tailored to the application it is linked against, including protection from speculative read gadgets, with low overhead.

Industry work is moving in the same direction. Trail of Bits proposed constant-time selection support for LLVM (\code{llvm.ct.select}, pull request \#166702) so that a compiler would be obliged to keep a selection branch-free~\cite{tobctselect}. Development continued through 2026. Rust is listed as a possible future user, but no Rust integration was found at the time of writing. An earlier Trail of Bits series examined how Rust cryptographic code relies on LLVM optimisation barriers~\cite{tobrust}.

\section{Constant-time analysis tools}
\label{sec:lit_tools}

Several tools detect constant-time violations without timing measurements. \textbf{ctgrind}~\cite{ctgrind} marks secret memory as uninitialised for Valgrind's memcheck, so any branch or address that depends on a secret is reported. \textbf{MicroWalk}~\cite{microwalk} (ACSAC 2018) combines dynamic binary instrumentation with mutual-information analysis to find and quantify memory- and control-flow-based leakage, and was applied to Intel IPP and Microsoft CNG. \textbf{haybale-pitchfork}~\cite{pitchfork} checks constant-time properties of LLVM IR by symbolic execution, which makes it directly applicable to Rust, but it documents support for old LLVM versions only. \textbf{AutoBench-CT}~\cite{autobenchct} (IACR ePrint 2026/2682) is a benchmark of 64 cryptographic tasks with manually verified ground truth, used to evaluate seven constant-time analysis tools; the authors report that ctgrind remains the most widely used tool in practice.

\section{The Rust cryptography ecosystem}
\label{sec:lit_rust}

Rust cryptographic libraries such as \code{curve25519-dalek} and the RustCrypto crates rely on the \code{subtle} crate~\cite{subtle} for constant-time selection and comparison. \code{subtle} represents secret booleans as \code{Choice}, an opaque byte, and by default passes it through a volatile read so that the compiler cannot treat it as a boolean and introduce a branch. Since version 2.5 an optional feature uses \code{core::hint::black\_box} instead. The Rust standard library documents \code{black\_box} as best-effort and not suitable for security guarantees~\cite{blackbox}.

Two real vulnerabilities frame the problem. RUSTSEC-2025-0144 (CVE-2026-22705)~\cite{rustsec2025_0144} is a timing side channel in the RustCrypto \code{ml-dsa} crate: the Decompose step used a variable-time integer division in the source code. It was fixed by replacing the division with a constant-time Barrett reduction. It is a source-level bug, not a compiler-introduced one, and is useful as a known positive control. CVE-2025-12888~\cite{cve2025_12888} describes a timing side channel in an X25519 implementation on Xtensa (ESP32) processors attributed to compiler optimisation and architecture limitations, a real-world example of the compiler-induced class.

\section{Comparative summary and research gap}
\label{sec:lit_gap}

Table~\ref{tab:lit} summarises the reviewed work.

\begin{table}[htbp]
\centering
\small
\renewcommand{\arraystretch}{1.18}
\begin{tabular}{@{}P{0.15\textwidth}P{0.26\textwidth}P{0.51\textwidth}@{}}
\toprule
Ref. & Area & What remains open for this project \\
\midrule
\cite{kocher96,bernstein05,osvik06} & Timing and cache attacks & Behaviour on 128-byte-line, large-L1 Apple cores \\
\cite{tvla,dudect,welch47} & Leakage assessment & Need for controls and bounds on non-detection \\
\cite{kasper09,hamburg09} & Constant-time AES & Empirical confirmation on a specific processor \\
\cite{gofetch,slap,flop} & Apple microarchitecture & Effect of predictors on leakage tests is unquantified \\
\cite{breakingbad,funwithflags,docc} & Compiler-induced violations & C libraries; Rust-specific mechanisms not studied \\
\cite{robocop,tobctselect} & Compiler defences & Rust exposure of \code{ct.select} not available \\
\cite{ctgrind,microwalk,pitchfork,autobenchct} & Analysis tools & Static/dynamic findings rarely linked to measured timing \\
\cite{subtle,blackbox,rustsec2025_0144} & Rust ecosystem & No empirical study of barrier survival across rustc settings \\
\bottomrule
\end{tabular}
\caption{Summary of the reviewed literature and the open points.}
\label{tab:lit}
\end{table}

Two gaps follow. The first, addressed in Part I, is practical: a course-scale, fully controlled measurement of how classic AES implementation strategies behave on a current Apple processor, with every claim checked against the assembly. The second, addressed by the plan in Part II, is a research gap. Compiler-induced violations have been studied at scale for C, but the constant-time mechanisms specific to Rust (\code{subtle}'s \code{Choice} and its barriers, and the boundary where a \code{Choice} is converted back into a \code{bool}) and Rust-specific build settings (MIR optimisation, monomorphisation, codegen units, LTO modes) have not been studied empirically. Neither has the question of which binary-level violations actually become measurable at runtime.
